\chapter{The Type System}\label{ch04:chapter}

\section{Design stance}\label{ch04:sec:stance}

Every later chapter of this monograph --- the board model, the design-rule
checker, the simulator, the geometry profiles, the TKS adapter, and the
adverse source audit --- rests on a single decision made at the lowest layer
of the candidate: \emph{semantic distinctions are enforced by types, not by
convention}. The CAS PCBA~8 type system is the place where that decision is
made concrete. It fixes, once, which values may exist at all; everything
downstream can then reason about which values may be \emph{connected},
\emph{checked}, \emph{simulated}, or \emph{adapted}, without re-litigating
what the values are.

The stance has four components, each of which recurs throughout this
chapter.

\begin{ruleenv}[Exact-type discipline]\label{ch04:rule:exact}
Every typed record in the candidate accepts only its exact runtime type.
Subclasses are rejected at construction and again at every revalidation.
A value that merely \emph{behaves like} a state literal, a port
specification, or a calibration reference is not one; the candidate refuses
to consume it.
\end{ruleenv}

\begin{ruleenv}[Fail-closed construction]\label{ch04:rule:failclosed}
Construction validates; revalidation re-validates. A record that cannot
prove its own well-formedness at the moment it is touched raises a typed
error rather than degrading into a permissive default. No field is ever
filled in silently: an absent Drive is absent, an unknown unit is an error,
an unsupported value type has no canonical identity and therefore cannot
enter a trace.
\end{ruleenv}

\begin{ruleenv}[Identity from structure]\label{ch04:rule:identity}
Wherever a value participates in evidence --- board digests, trace events,
measurement records, geometry profiles --- its identity is a canonical,
versioned, deterministic serialization of its complete typed structure,
hashed with SHA-256. Two values are the same exactly when their canonical
payloads are byte-identical. There is no ambient notion of equality that
can be substituted for this one.
\end{ruleenv}

\begin{ruleenv}[Boundary in the types]\label{ch04:rule:boundary}
The research boundary --- Ring~4, \textsc{open}, \textsc{not licensed},
\textsc{held} --- is itself a type. Artifacts do not carry a mutable status
field that a caller could upgrade; they carry a value of a type whose only
inhabitant \emph{is} the fixed boundary, rederived and rechecked by every
public operation.
\end{ruleenv}

Nothing in this chapter is a claim about cognition, about any deployed
model, or about physical electronics. The type system is a software
contract. Where a type carries a CAS name (Noetic, Domain, Drive,
Register), the name records which CAS coordinate the type is \emph{intended}
to represent inside the research candidate; it does not assert that the
coordinate is realized in any external system. That boundary is restated
formally in Section~\ref{ch04:sec:boundary}.

\section{The four semantic factors}\label{ch04:sec:factors}

The candidate models the CAS state space as a product of finite coordinate
sets. Four enumerated factor types are declared, and --- deliberately ---
nothing else is assumed about their interaction. In particular, the
implementation's own module documentation is explicit that the factors are
\emph{not} called geometrically orthogonal: no orthogonality theorem is
assumed by the type layer, and none is imported from the source book.

\begin{definition}[Noetic factor]\label{ch04:def:noetic}
$\mathrm{Noetic}$ is a ten-element enumeration carried on the integers
$0$ through $9$: \textsc{idea}~(0), \textsc{mind}~(1),
\textsc{positive}~(2), \textsc{negative}~(3), \textsc{vibration}~(4),
\textsc{female}~(5), \textsc{male}~(6), \textsc{rhythm}~(7),
\textsc{above-cause}~(8), \textsc{below-effect}~(9).
\end{definition}

\begin{definition}[Domain factor]\label{ch04:def:domain}
$\mathrm{Domain}$ is a four-element enumeration carried on the strings
\texttt{P}, \texttt{C}, \texttt{E}, \texttt{O} (Principled, Cognitive,
Evaluative, Operational). Each domain exposes a total precedence rank
$P \mapsto 0$, $C \mapsto 1$, $E \mapsto 2$, $O \mapsto 3$, encoding the
structural order $P \succ C \succ E \succ O$ as data available to the
checker; the type layer itself attaches no descent semantics to the rank.
\end{definition}

\begin{definition}[Drive factor]\label{ch04:def:drive}
$\mathrm{Drive}$ is a seven-element enumeration carried on the strings
\texttt{F1} through \texttt{F7}.
\end{definition}

\begin{definition}[Register factor]\label{ch04:def:register}
$\mathrm{Register}$ is a four-element enumeration carried on the strings
\texttt{T} (Transcendent), \texttt{A} (Analytical), \texttt{f} (Affective),
and \texttt{E} (Embodied). The lowercase \texttt{f} is deliberate and
canonical; it prevents a case-folding collision with the Embodied register
\texttt{E} and with the Evaluative domain \texttt{E}, which lives in a
different enumeration entirely.
\end{definition}

\begin{remark}[Domain \texttt{E} versus Register \texttt{E}]
The string \texttt{E} appears as a carrier value in both
$\mathrm{Domain}$ and $\mathrm{Register}$. Under the exact-type discipline
of Rule~\ref{ch04:rule:exact} this is harmless: the two values inhabit
disjoint enumeration types and no operation in the candidate accepts a bare
string where a factor is required. The collision would be hazardous only in
a stringly-typed design --- which is precisely the design the candidate
refuses. This is a small, early illustration of why the type system exists.
\end{remark}

\section{State literals}\label{ch04:sec:literals}

\begin{definition}[State literal]\label{ch04:def:stateliteral}
A state literal is an element of
\[
\mathrm{StateLiteral}
 = \mathrm{Noetic} \times \mathrm{Domain}
 \times \mathrm{Drive}_{\mathrm{opt}} \times \mathrm{Register}_{\mathrm{opt}},
\]
where $\mathrm{Drive}_{\mathrm{opt}} = \mathrm{Drive} \sqcup \{\bot\}$ and
$\mathrm{Register}_{\mathrm{opt}} = \mathrm{Register} \sqcup \{\bot\}$
adjoin an explicit absence marker. In the implementation the type is the
frozen record \texttt{CasState} with fields \texttt{noetic},
\texttt{domain}, \texttt{drive} (optional), and \texttt{register}
(optional).
\end{definition}

Three points about this definition carry analytical weight.

First, \emph{optionality is explicit, not defaulted}. The implementation's
own docstring states the contract: Drive and Register are optional because
the inherited notation permits partial literals, and ``their absence is
explicit and never filled by a default or inferred from another state.''
A literal with no Drive is a different mathematical object from a literal
with any particular Drive, and the candidate never manufactures the latter
from the former. This matters because the source notation
(\texttt{\^{}n\^{}D}, \texttt{\^{}n\^{}D\_\{Fk\}},
\texttt{\^{}n\^{}D\_\{Fk\}|R}) treats the shorter forms as legitimate
partial specifications, not as abbreviations of some canonical completion.
A type system that auto-completed them would be inventing semantics the
source does not license.

Second, \emph{construction and revalidation enforce factor exactness}. Each
field must be an exact member of its factor enumeration; \texttt{None} is
admitted only for the two optional fields. A string \texttt{"P"} in the
domain position, an integer \texttt{3} in the noetic position, or a
subclass of any factor type is rejected with a typed error.

Third, \emph{the literal has a canonical rendering}. The record exposes a
rendering of the form \texttt{\^{}3\^{}P}, \texttt{\^{}3\^{}P\_\{F2\}}, or
\texttt{\^{}3\^{}P\_\{F2\}|A}, produced only after revalidation. The
rendering is derived output, never an alternative identity: equality and
hashing are structural, over the typed fields, and the canonical trace
identity of a state literal (Section~\ref{ch04:sec:traceid}) serializes the
fields individually rather than the rendered string. A rendered string can
be regenerated from the fields; the fields cannot be safely regenerated
from an arbitrary string without a parser, and the v0.1 candidate
deliberately declares no parser (Chapter~10 records why).

\section{ND-cells and the forgetful projection}\label{ch04:sec:ndcells}

\begin{definition}[ND-cell]\label{ch04:def:ndcell}
An ND-cell is an element of
\[
\mathrm{NDCell} = \mathrm{Noetic} \times \mathrm{Domain},
\]
implemented as the frozen record \texttt{NdCell} with typed fields
\texttt{domain} and \texttt{noetic} and a derived label of the form
\texttt{P0}, \texttt{C7}, \texttt{O9}. The forty canonical cells are fixed
in the order $P0,\ldots,P9,\,C0,\ldots,C9,\,E0,\ldots,E9,\,O0,\ldots,O9$,
exported as an immutable tuple together with the parallel tuple of canonical
labels.
\end{definition}

The canonical order is load-bearing: it is the coordinate order of every
simplex address (Section~\ref{ch04:sec:simplex}), the row discipline of
every geometry identity payload, and the order against which the mixed-simplex
validator checks incoming cell tuples. A permutation of the forty cells is
not ``the same cells in a different order''; it is a malformed address.
Table~\ref{ch04:tab:cells} fixes the correspondence between coordinate
index, label, and factor pair once, for reference throughout the
monograph.

\begin{longtable}{@{}rllcrll@{}}
\caption{The forty canonical ND-cells in canonical order. Index $k$ is the
coordinate position used by every simplex address and geometry payload.}
\label{ch04:tab:cells}\\
\toprule
$k$ & Label & Factors & \hspace{2em} & $k$ & Label & Factors \\
\midrule
\endfirsthead
\toprule
$k$ & Label & Factors & \hspace{2em} & $k$ & Label & Factors \\
\midrule
\endhead
\bottomrule
\endfoot
1 & \texttt{P0} & (Idea, Principled) & & 21 & \texttt{E0} & (Idea, Evaluative) \\
2 & \texttt{P1} & (Mind, Principled) & & 22 & \texttt{E1} & (Mind, Evaluative) \\
3 & \texttt{P2} & (Positive, Principled) & & 23 & \texttt{E2} & (Positive, Evaluative) \\
4 & \texttt{P3} & (Negative, Principled) & & 24 & \texttt{E3} & (Negative, Evaluative) \\
5 & \texttt{P4} & (Vibration, Principled) & & 25 & \texttt{E4} & (Vibration, Evaluative) \\
6 & \texttt{P5} & (Female, Principled) & & 26 & \texttt{E5} & (Female, Evaluative) \\
7 & \texttt{P6} & (Male, Principled) & & 27 & \texttt{E6} & (Male, Evaluative) \\
8 & \texttt{P7} & (Rhythm, Principled) & & 28 & \texttt{E7} & (Rhythm, Evaluative) \\
9 & \texttt{P8} & (Above/Cause, Principled) & & 29 & \texttt{E8} & (Above/Cause, Evaluative) \\
10 & \texttt{P9} & (Below/Effect, Principled) & & 30 & \texttt{E9} & (Below/Effect, Evaluative) \\
11 & \texttt{C0} & (Idea, Cognitive) & & 31 & \texttt{O0} & (Idea, Operational) \\
12 & \texttt{C1} & (Mind, Cognitive) & & 32 & \texttt{O1} & (Mind, Operational) \\
13 & \texttt{C2} & (Positive, Cognitive) & & 33 & \texttt{O2} & (Positive, Operational) \\
14 & \texttt{C3} & (Negative, Cognitive) & & 34 & \texttt{O3} & (Negative, Operational) \\
15 & \texttt{C4} & (Vibration, Cognitive) & & 35 & \texttt{O4} & (Vibration, Operational) \\
16 & \texttt{C5} & (Female, Cognitive) & & 36 & \texttt{O5} & (Female, Operational) \\
17 & \texttt{C6} & (Male, Cognitive) & & 37 & \texttt{O6} & (Male, Operational) \\
18 & \texttt{C7} & (Rhythm, Cognitive) & & 38 & \texttt{O7} & (Rhythm, Operational) \\
19 & \texttt{C8} & (Above/Cause, Cognitive) & & 39 & \texttt{O8} & (Above/Cause, Operational) \\
20 & \texttt{C9} & (Below/Effect, Cognitive) & & 40 & \texttt{O9} & (Below/Effect, Operational) \\
\end{longtable}

\begin{definition}[Forgetful projection]\label{ch04:def:projection}
The projection
\[
\pi : \mathrm{StateLiteral} \longrightarrow \mathrm{NDCell},
\qquad
\pi(n, d, f, r) = (n, d)
\]
forgets the Drive and Register coordinates. In the implementation it is the
method \texttt{to\_nd\_cell}, whose docstring states that it projects
``explicitly forgetting two axes.''
\end{definition}

\begin{proposition}[The projection is many-to-one and not an isomorphism]
\label{ch04:prop:manytoone}
Each fiber $\pi^{-1}(c)$ over a canonical cell $c$ has cardinality
$(7+1)\cdot(4+1) = 40$, so
\[
\lvert \mathrm{StateLiteral} \rvert = 40 \cdot 40 = 1600,
\qquad
\lvert \mathrm{NDCell} \rvert = 40 .
\]
In particular $\pi$ is surjective but not injective, and no map
$s : \mathrm{NDCell} \to \mathrm{StateLiteral}$ satisfies
$s \circ \pi = \id_{\mathrm{StateLiteral}}$.
\end{proposition}

\begin{proof}
Fix a cell $c = (n, d)$. A literal in the fiber is determined by its Drive
coordinate, ranging over the seven drives together with the explicit
absence $\bot$, and its Register coordinate, ranging over the four
registers together with $\bot$: $8 \cdot 5 = 40$ choices, all realized and
pairwise distinct by construction of the record type. Surjectivity is
witnessed by the partial literal $(n, d, \bot, \bot)$. Injectivity fails
because any fiber contains more than one element; a left inverse
$s$ with $s \circ \pi = \id$ would force injectivity of $\pi$.
\end{proof}

\subsection{Why forgetting Drive and Register matters analytically}
\label{ch04:sec:whyforget}

Proposition~\ref{ch04:prop:manytoone} is elementary as counting; its
consequences are not. The candidate keeps both types --- and keeps the
projection explicit --- because three distinct analytical hazards arise
when the collapse is performed silently.

\emph{First hazard: aggregation is a modeling decision, not a
formality.} A simplex address (next section) distributes mass over the
forty ND-cells. When a state literal carrying a Drive and a Register is
scored into that address, forty distinct literals land in the same
coordinate. Whatever distinguished analytical work the Drive coordinate was
doing --- which motivational architecture the state participates in --- and
whatever the Register coordinate was doing --- which operational character
the state manifests under --- is unrecoverable from the address. Under
Correction~2 of the framework's own mathematical discipline (prerequisite
satisfaction is one-directional; a retraction is not an inverse), this is
exactly the situation of a section-retraction failure: one can embed a cell
into a canonical partial literal and project back to the same cell, but one
cannot project a full literal to a cell and re-embed without loss. The type
system makes the loss visible as a typed operation
(\texttt{to\_nd\_cell}) rather than an implicit coercion, so that any
analysis pipeline that crossed the projection can be audited for whether it
was entitled to.

\emph{Second hazard: distinctness claims silently weaken.} Two literals
$\ell_1 = ({\wedge}3, P, F2, A)$ and $\ell_2 = ({\wedge}3, P, F6, f)$ are
distinct objects with, on the framework's own account, different
diagnostic content. Their projections coincide. A report that compared
board state bindings at cell resolution and found ``the same state on both
ports'' would be making a claim that is true of the projections and false
of the literals. Because the candidate's canonical trace identity for a
state value serializes all four fields (Section~\ref{ch04:sec:traceid}),
identity-level evidence never passes through the projection; only an
analysis that explicitly chose cell resolution can conflate the two, and
that choice is then legible in the analysis, not buried in the types.

\emph{Third hazard: the inverse direction invites fabrication.} Since no
section of $\pi$ exists (Proposition~\ref{ch04:prop:manytoone}), any
procedure that produces a full literal from cell-resolution data has
\emph{added} information --- a Drive attribution and a Register attribution
--- that the data did not contain. In an evidence-bearing pipeline that
addition must be declared as a modeling step with its own provenance, or it
is fabrication. The type system cannot prevent an analyst from writing such
a procedure, but it can and does refuse to provide one: there is no
``default drive,'' no ``neutral register,'' and no constructor that lifts a
cell to a literal by fiat. The only canonical object available at cell
resolution is the cell itself.

\begin{remark}[Relation to the register-multiplicity discipline]
The v7.7 framework layer distinguishes register-set specification,
co-presence, and cross-register integration precisely because collapsing
them destroys diagnostic content. The type layer of this candidate sits
below that discipline --- a single \texttt{CasState} carries at most one
register --- but the same conservative instinct governs both: where the
notation distinguishes, the types distinguish, and any collapse is an
explicit, auditable act. Multi-register expressions are the business of the
notation and its (future) parser, not of the state-literal record; the
record's single optional register is the correct target for one
register-specified primitive, not a claim that richer expressions do not
exist.
\end{remark}

\section{Simplex addresses}\label{ch04:sec:simplex}

\begin{definition}[SimplexAddress40]\label{ch04:def:simplex}
A $\mathrm{SimplexAddress}_{40}$ is a strictly positive normalized profile
over the forty canonical ND-cells:
\[
p \in \mathbb{R}^{40},
\qquad p_k > 0 \ \text{for all } k,
\qquad \sum_{k=1}^{40} p_k = 1 .
\]
Equivalently, $p$ is a point of the \emph{open} 39-dimensional probability
simplex
\[
\Delta_{39}^{\circ}
 = \Bigl\{\, p \in \mathbb{R}^{40} \ :\ p_k > 0,\ \textstyle\sum_k p_k = 1 \,\Bigr\},
\]
with coordinates indexed in the canonical cell order of
Definition~\ref{ch04:def:ndcell}. The implementation is the frozen record
\texttt{MixedSimplexAddress} with exactly forty finite, strictly positive
\texttt{Decimal} weights whose \emph{exact} sum is $1$, cells fixed to the
canonical tuple, and a fixed profile discriminator
\texttt{CANONICAL\_MIXED\_SIMPLEX\_BASELINE} under the versioned profile
name \texttt{cas-pcba-v8.mixed-simplex.v1}.
\end{definition}

\subsection{Why strict positivity}\label{ch04:sec:positivity}

The open simplex --- rather than the closed simplex
$\overline{\Delta}_{39}$, which admits zero coordinates --- is a deliberate
and threefold choice.

\emph{Geometric reason.} The Fisher--Rao metric, which
Chapter~\ref{ch06:chapter} develops as the baseline geometry, has
coordinate expression
\[
g_p(u, v) = \sum_{k} \frac{u_k v_k}{p_k}
\]
on the tangent space at $p$. The expression degenerates as any
$p_k \to 0$: the metric coefficients blow up, the closed simplex is not a
Riemannian manifold under this metric, and geodesics can reach the boundary
only in a limiting sense. Restricting the type to the open simplex is what
makes ``the geometry profile'' a coherent phrase at all: every value the
type admits is a point where the metric is defined and smooth
\cite{amari,ay}. The type system thus enforces, at construction time, the
precondition of every downstream geometric computation, instead of letting
each computation rediscover boundary pathologies at run time.

\emph{Semantic reason.} A zero coordinate is not a small quantity; it is a
categorical assertion --- \emph{this cell carries no mass whatsoever}. That
is a different and much stronger claim type than a profile that
concentrates mass elsewhere, and it is a claim the candidate's evidence
discipline has no basis to admit: no measurement in the admitted evidence
distinguishes ``vanishingly small'' from ``structurally absent.'' By
excluding zeros the type refuses to let an address smuggle in an
impossibility claim. Analysts who need absence claims must make them in a
different vocabulary (the framework's lacunarity analysis), where they can
carry their own evidence burden.

\emph{Numerical reason.} Downstream computations divide by coordinates and
take logarithms and square roots of them. Strict positivity, validated
strictly at construction and at every revalidation, means no downstream
profile computation needs a guard against division by zero on data that
passed the type. The validation is exact, not tolerance-based: weights are
\texttt{Decimal} values, positivity is exact comparison, and the sum
constraint is checked with an exact aligned-integer summation
(Section~\ref{ch04:sec:decimal}) rather than a floating tolerance. Distinct
questions --- ``is this a well-formed address?'' and ``are these two
addresses numerically distinguishable?'' --- get distinct answers from
distinct machinery; the second is the pseudometric disclosure of
Chapter~\ref{ch06:chapter}.

\subsection{Exact normalization}\label{ch04:sec:normalization}

The type provides one constructor from unnormalized data:
\texttt{from\_positive\_weights} accepts exactly forty finite, strictly
positive \texttt{Decimal} values, divides by their exact sum under the
fixed decimal context, and then repairs the residual: the exact sum of the
normalized weights is recomputed, the deviation from exactly $1$ is
computed by exact arithmetic, and that residual is added to a
deterministically chosen anchor coordinate (the largest input weight, ties
broken by lowest index). If the repair would drive the anchor coordinate
nonpositive, construction fails rather than admitting a boundary point.
The result is an address whose weights sum to \texttt{Decimal("1")}
\emph{exactly} --- not to within an epsilon --- so the sum constraint in
Definition~\ref{ch04:def:simplex} is a checkable identity, not an
approximation, and the constructor is deterministic: the same input tuple
yields the same address and the same identity digest on every run.

\subsection{What a simplex address is not}\label{ch04:sec:simplexnot}

\begin{ruleenv}[Not one-hot]\label{ch04:rule:notonehot}
A one-hot vector --- mass $1$ on one cell, $0$ elsewhere --- is a vertex of
the closed simplex and is excluded from $\Delta_{39}^{\circ}$ by strict
positivity. A simplex address therefore can never assert ``the system is in
cell $c$.'' The object that names a single cell is \texttt{NdCell} itself;
the object that names a single full state is \texttt{CasState}. Conflating
a sharply concentrated address with a definite cell is a category error the
type system makes unrepresentable rather than merely inadvisable.
\end{ruleenv}

\begin{ruleenv}[Not a finite TKS kernel element]\label{ch04:rule:nottks}
The TKS~v8 finite reference kernel traffics in finite carriers: tagged
finite values, finite transition tables, finite arrows. A simplex address
has an uncountable ambient type and continuous structure; it is not an
element of any finite carrier, and the adapter boundary
(Chapter~9, and the boundary discipline restated in the
seed specification) rejects any attempt to coerce one into a finite arrow.
The signal-kind vocabulary of Section~\ref{ch04:sec:ports} keeps the two
worlds apart at the wiring level: \texttt{MIXED\_SIMPLEX\_ADDRESS} and
\texttt{FINITE\_STATE} are distinct signal kinds with disjoint admissible
units, so a net cannot silently deliver a continuous profile to a port that
expects a finite state.
\end{ruleenv}

\begin{remark}[Identity of an address]
Every address exposes a canonical identity payload --- profile name,
profile kind, the decimal-context digest, the forty canonical labels in
order, and the forty weights as tagged decimal strings --- and a SHA-256
over that payload. Two addresses are evidentially the same exactly when
these digests agree. This is the exact-identity discipline that
Chapter~\ref{ch06:chapter} contrasts with the binary64 pseudometric: the
digest separates all distinct addresses; the floating-point distance does
not promise to.
\end{remark}

\section{Ports: direction, signal kind, unit}\label{ch04:sec:ports}

Ports are where the type system meets wiring. A port specification is a
frozen record with five typed fields.

\begin{definition}[Port specification]\label{ch04:def:portspec}
A \texttt{PortSpec} consists of: a reference string \texttt{ref} conforming
to the safe-reference profile (Section~\ref{ch04:sec:provenance}); a
\texttt{direction} in the two-element enumeration
$\{\textsc{input}, \textsc{output}\}$; a \texttt{signal\_kind} in a
nine-element enumeration; a \texttt{unit} in a nine-element enumeration;
and a boolean \texttt{required} (default true).
\end{definition}

The signal kinds are: \textsc{activation}, \textsc{decision},
\textsc{cas\_state}, \textsc{finite\_state},
\textsc{mixed\_simplex\_address}, \textsc{projector\_profile},
\textsc{text}, \textsc{clock}, and \textsc{scalar}. The units are:
\textsc{none}, \textsc{activation\_0\_1}, \textsc{boolean},
\textsc{probability}, \textsc{normalized\_0\_10},
\textsc{cas\_ohm\_score}, \textsc{volt}, \textsc{second}, and
\textsc{hertz}.

\begin{ruleenv}[Closed signal-kind/unit admissibility]\label{ch04:rule:sigunit}
Admissible pairings of signal kind and unit form a closed, immutable table
(Table~\ref{ch04:tab:sigunit}). The admissibility predicate returns
false --- it does not raise, and it does not guess --- for any pair
involving a value that is not an exact member of the respective
enumeration.
\end{ruleenv}

\begin{longtable}{@{}lp{7.2cm}@{}}
\caption{Closed signal-kind/unit admissibility table.}
\label{ch04:tab:sigunit}\\
\toprule
Signal kind & Admissible units \\
\midrule
\endfirsthead
\toprule
Signal kind & Admissible units \\
\midrule
\endhead
\bottomrule
\endfoot
\textsc{activation} & \textsc{activation\_0\_1} \\
\textsc{decision} & \textsc{boolean} \\
\textsc{cas\_state} & \textsc{none} \\
\textsc{finite\_state} & \textsc{none} \\
\textsc{mixed\_simplex\_address} & \textsc{probability} \\
\textsc{projector\_profile} & \textsc{none} \\
\textsc{text} & \textsc{none} \\
\textsc{clock} & \textsc{hertz} \\
\textsc{scalar} & \textsc{none}, \textsc{probability},
\textsc{normalized\_0\_10}, \textsc{cas\_ohm\_score}, \textsc{volt},
\textsc{second}, \textsc{hertz} \\
\end{longtable}

The table encodes a semantic commitment per row. An activation is a
bounded intensity in $[0,1]$ and nothing else; a decision is a boolean and
nothing else; the four ``carrier-typed'' kinds (CAS state, finite state,
projector profile, text) are unitless because their values are structured
identities rather than magnitudes; a mixed-simplex address is
probability-typed because its coordinates are exactly a normalized
profile; a clock is a rate. Only the deliberately generic \textsc{scalar}
kind admits a menu of units, and it is correspondingly excluded from every
semantically loaded role: no reference component contract in the closed
catalog exposes a scalar port where an activation is meant, so a scalar
can never impersonate an activation merely by carrying a compatible
number.

Two remarks on this table. First, the presence of \textsc{volt} among the
scalar units is a naming admission, not a physics claim: the unit exists so
that source-derived example boards can be transcribed faithfully, and the
research boundary (Section~\ref{ch04:sec:boundary}) denies physical
telemetry regardless of unit names. A scalar labeled volts is a scalar
labeled volts; nothing in the candidate measures electric potential.
Second, the table is deliberately closed: an unsupported pairing is not a
warning but a structural impossibility that the DRC reports under its net
diagnostics. Extending the table is a versioned change to the engine
identity, visible in every downstream digest.

\section{Component kinds and typed parameter contracts}
\label{ch04:sec:contracts}

A component on a board names a \emph{kind}; the kind resolves, through a
closed catalog, to a \emph{contract} that fixes what the component is
allowed to look like.

\begin{definition}[Parameter rule]\label{ch04:def:paramrule}
A \texttt{ParameterRule} declares, for one named parameter: a value type
(\textsc{decimal} or \textsc{string}); whether the parameter is required;
optional finite \texttt{Decimal} bounds \texttt{minimum} and
\texttt{maximum} (with an exclusive-maximum flag); and, for string
parameters, a closed set of admissible choices, each of which must itself
conform to the safe-reference profile. Construction rejects a minimum
exceeding a maximum, non-finite bounds, and untyped flags.
\end{definition}

\begin{definition}[Component contract]\label{ch04:def:contract}
A \texttt{ComponentContract} binds: a kind name and a versioned semantic
profile string; a tuple of \texttt{PortSpec} values; an immutable mapping
from parameter names to \texttt{ParameterRule} values; a set of required
state-binding names; and three booleans --- \texttt{deterministic},
\texttt{stateful}, and \texttt{tks\_finite\_adapter\_eligible} --- which
downstream stages consult before simulating or adapting the component.
\end{definition}

The default catalog in the v0.1 implementation is small and closed. Its
kinds are: \texttt{pass} (identity on one activation), \texttt{scale}
(activation scaling with a nonnegative decimal \texttt{factor} and an
activation-level \texttt{ceiling}), \texttt{cjt} (the cognitive-junction
reference: two activation inputs \texttt{cause} and \texttt{bias}, one
activation output \texttt{effect}; parameters \texttt{beta} (nonnegative
finite decimal), \texttt{saturation}, and \texttt{cutoff} (both
activation-level); required state bindings \texttt{cause},
\texttt{bias}, \texttt{effect}), \texttt{decision\_threshold}
(activation in, boolean decision out, with an activation-level
\texttt{threshold}), \texttt{mux} (two activation inputs selected by a
decision input), \texttt{fuse} (activation in and out, with an
activation-level \texttt{threshold} and a string \texttt{action}
restricted to the closed choice set \texttt{permanent\_block} or
\texttt{warn}; stateful, because a permanently blocking fuse carries
run-local open/closed state), \texttt{delay} (the unit delay, stateful,
with an activation-level \texttt{initial}), and \texttt{finite\_transition}
(finite state in and out, the only kind eligible for the TKS finite
adapter). The catalog as a whole has a canonical digest that enters the
engine identity and hence every board identity
(Chapter~\ref{ch05:chapter}); adding a kind, changing a bound, or widening
a choice set is a visible engineering event, not a quiet edit.

Two consequences deserve emphasis. The \texttt{beta} parameter of the
reference junction is typed as a \emph{finite} nonnegative decimal: the
source book's idealized infinite gain (the $\beta = \infty$ of the
published Cognitive Safety Board, Chapter~11) is
unrepresentable in the candidate's parameter types, which converts a
documented source inconsistency into a construction-time impossibility.
And the \texttt{fuse} kind's \texttt{action} parameter is a closed string
choice, not free text: there is no way to declare a fuse whose action the
simulator would have to guess at.

\section{Nets, endpoints, and driver/consumer discipline}
\label{ch04:sec:nets}

\begin{definition}[Endpoint and net]\label{ch04:def:net}
An \texttt{Endpoint} is a pair of reference strings
(\texttt{owner\_ref}, \texttt{port\_ref}); the distinguished owner sentinel
\texttt{\$board} denotes the board's own external port surface and is the
only reference exempt from the leading-alphanumeric rule of the
safe-reference profile. A \texttt{Net} is a reference string, a tuple of
endpoints, and a boolean \texttt{safety\_critical} flag.
\end{definition}

At the type layer a net is deliberately dumb: it is a set of typed
attachment points, nothing more. The intelligence lives in the design-rule
checker (Chapter~7), which resolves every endpoint to a
declared owner and port, requires exactly one driver per net, and enforces
signal-kind and unit agreement between the driver and every consumer under
the closed table of Rule~\ref{ch04:rule:sigunit}. The type layer's
contribution is to make the checker's job well-posed: because endpoints are
exact typed records with validated references, the checker's questions
(``does this owner exist?'', ``does this port exist on that owner?'',
``what kind and unit does it carry?'') have deterministic answers, and
because nets are immutable, the answers cannot change between the check and
the use of its report. The driver/consumer rules themselves ---
\emph{one} driving output per net, at least one load, direction
compatibility, no multiply-connected port, required inputs connected or
defaulted in the selected simulator profile --- are stated and itemized as
\texttt{NET}-class diagnostics in Chapter~7; here they are
named only to record that every rule is expressible as a predicate over the
exact types of this chapter, with no reflection, no string parsing, and no
appeal to convention.

\section{Test points and calibration references}\label{ch04:sec:tpcal}

\begin{definition}[Test point]\label{ch04:def:testpoint}
A \texttt{TestPoint} names a reference, the net it observes
(\texttt{net\_ref}), an optional \texttt{calibration\_ref}, and a boolean
\texttt{safety\_critical} flag.
\end{definition}

\begin{definition}[Calibration reference]\label{ch04:def:calref}
A \texttt{CalibrationRef} carries: reference and version strings; a
\texttt{profile\_sha256} in exact uppercase 64-hex form; a typed
\texttt{measurement\_kind} drawn from the six-element enumeration
\textsc{design\_value}, \textsc{behavioral\_proxy},
\textsc{instrumented}, \textsc{human\_score}, \textsc{llm\_score},
\textsc{model\_derived}; a typed \texttt{unit}; instrument and protocol
identity strings; a finite \texttt{Decimal} \texttt{uncertainty}; and a
boolean \texttt{packet\_eligible}, default false.
\end{definition}

The design intent is to make the \emph{shape} of measurement honesty a
type. A test point is a declared observation site, not an observation; a
calibration reference is a declared pointer to calibration metadata, not
evidence that calibration occurred. The fields force every calibration
claim to name its instrument identity, protocol identity, uncertainty, and
measurement kind at declaration time --- so the later question ``what would
it take for this observation to count?'' has a typed answer --- while the
\texttt{packet\_eligible} flag records the candidate's standing rule that
locally constructed calibration records are structural, test-only objects:
nothing a board author writes can make a calibration reference
Evidence-Packet eligible, because eligibility is granted by an
authentication process outside the type system, never by the record
itself. The measurement-kind enumeration also encodes one substantive
prohibition enforced at simulation intake: a board \emph{input} sample may
not claim \textsc{model\_derived} provenance, which blocks the circular
move of feeding a model-derived quantity back in as if it were fresh
external data.

\section{Deterministic trace and measurement identity}
\label{ch04:sec:traceid}

Everything the simulator emits is identity-bearing. The type layer
provides the vocabulary in three parts.

\emph{Canonical value identities.} Every value that can flow through a
port or appear in a trace has exactly one canonical tagged JSON identity
under the versioned profile \texttt{cas-pcba-v8.trace-value.v1}: finite
decimals (tagged with the decimal-context digest), tagged finite atoms
(Section~\ref{ch04:sec:tagged}), full CAS state literals (all four fields
serialized, including explicit nulls for absent Drive or Register),
mixed-simplex addresses (by their complete identity payload and its
digest), and projector profiles (likewise, with the matrix digest linked).
Any other runtime type has \emph{no} canonical identity and is rejected;
``unsupported value'' is an error, never a best-effort string. A canonical
identity round-trips: the serialized text must re-canonicalize to itself
byte-for-byte, and its digest is the SHA-256 of that canonical text. Port
validation is performed against the serialized identity, per signal kind:
an activation port's value must be a decimal identity in $[0,1]$; a
decision port's must be an exact boolean atom; a mixed-simplex port's must
be a mixed-simplex identity; and so on --- so a value's fitness for a port
is checked on the same bytes that enter the evidence, not on a live object
that could differ from what was recorded.

\emph{Measurement records.} A \texttt{MeasurementRecord} (profile
\texttt{cas-pcba-v8.measurement-record.v2}) binds: an identity-scope
digest tying the record to one exact run scope; a non-negative integer
tick; producer and port references; the canonical value text and its
digest; typed unit and measurement kind; an optional calibration
reference; the \texttt{packet\_eligible} flag, fixed false at creation;
an ordered tuple of \emph{child measurement identities}; and the fixed
research boundary. The record's own \texttt{measurement\_id} is the
SHA-256 of its complete canonical payload --- so a measurement cannot be
edited, re-parented, re-ticked, or re-valued without changing its
identity, and the child tuple gives each derived measurement an explicit,
acyclic ancestry that the result constructor validates (children must
exist; derivation must be acyclic; ticks and trace linkage must agree).
The discipline resembles the entity/activity/derivation separation of the
W3C provenance model \cite{prov}, and the candidate is explicit that the
resemblance is methodological only: no PROV-O conformance is claimed.

\emph{Trace events and results.} Trace events and simulation results carry
their own versioned profiles (\texttt{cas-pcba-v8.trace-event.v2},
\texttt{cas-pcba-v8.simulation-result.v2}) and bind the exact DRC report
identity, engine and simulator profiles, component catalog and
decimal-context digests, source-bundle digest, and research boundary, as
Chapter~8 details. What matters at the type layer is
uniformity: every one of those bindings is a typed field holding a
validated digest or a validated record, so ``the trace says X'' is always
reducible to ``these exact bytes, under this exact profile, hash to this
exact identity.''

\section{Research boundary and source provenance types}
\label{ch04:sec:provenance}

\begin{definition}[Research boundary type]\label{ch04:def:boundary}
\texttt{ResearchBoundary} is a class with \emph{no instance storage}
(empty slots). Its constructor validates that every argument equals the
fixed research posture --- admission ring $4$, status \textsc{open},
license status \textsc{not licensed}, disclosure status \textsc{held},
safety certification false, telemetry authorization false, regulatory
authority false, and the fixed research claim sentence --- and raises
otherwise. All accessors are constant properties. Equality holds exactly
between values of the exact type.
\end{definition}

The empty-slots design is the point: because instances store nothing, even
low-level attribute forcing cannot mutate a boundary value into an
``accepted,'' ``licensed,'' ``disclosed,'' or ``certified'' state --- there
is nowhere for the forged field to live. Every public operation that
consumes a boundary calls the revalidation function, which checks the exact
type and re-reads every property against the fixed posture. The boundary
travels inside measurement records, projector identities, and simulation
results, so the Ring-4 / \textsc{not licensed} posture is not a banner on
the paper; it is a field in the evidence.

\begin{definition}[Safe-reference profile]\label{ch04:def:saferef}
Every reference string in the candidate (component, port, net, test-point,
calibration, feedback, profile, basis, and source identifiers) must match
the versioned profile \texttt{cas-pcba-v8.safe-reference-ascii.v1}: ASCII
only, first character alphanumeric, up to 128 characters drawn from
alphanumerics, dot, underscore, colon, and hyphen. The validator rejects
control characters, bidirectional and formatting characters, and confusable
non-ASCII lookalikes.
\end{definition}

This is a human-review defense encoded as a type rule: an identifier that
will appear in audit reports, diffs, and evidence packets must not be able
to spoof another identifier under visual review or reorder text under
bidirectional rendering. The sole sanctioned exception is the
\texttt{\$board} endpoint sentinel of Definition~\ref{ch04:def:net}, which
is permitted only where an endpoint owner is expected and only by an
explicit flag at the call site.

\emph{Source provenance.} The candidate binds its own implementation
identity into everything it produces: a source-bundle digest over the
package's Python sources (profile
\texttt{cas-pcba-v8.python-source-bundle.v1}), an engine identity payload
(profile \texttt{cas-pcba-v8.reference-engine.v1}) that nests the source
identity, the component-catalog digest, and the complete decimal-context
payload with its digest. Board identities embed the engine identity;
traces embed both; so every artifact answers not only ``what was
modeled?'' but ``which exact software, under which exact arithmetic and
catalog, said so?''

\section{Tagged exact values}\label{ch04:sec:tagged}

\begin{definition}[Tagged finite value]\label{ch04:def:tagged}
A finite carrier atom is exactly one of: a boolean, an integer, or a
string. Its canonical identity is a tagged payload under the versioned
profile \texttt{cas-pcba-v8.tagged-finite-value.v1} carrying the type tag
(\texttt{bool}, \texttt{int}, or \texttt{string}) and a canonical value
representation --- booleans as JSON booleans, integers as their exact
decimal string, strings as themselves. Equality of finite values is
equality of tagged payloads: type tag \emph{and} value.
\end{definition}

\begin{ruleenv}[Type-pun exclusion]\label{ch04:rule:typepun}
Integer $1$, Boolean \emph{true}, and string \texttt{"true"} are three
distinct values everywhere in the candidate: in DRC evaluation, in
execution, in traces, in measurement identities, and across the TKS
adapter. No implicit conversion between them exists at any layer.
\end{ruleenv}

The hazards this rule blocks are concrete, and each corresponds to a
failure mode observed in mainstream ecosystems.

\emph{The bool/int pun.} In the implementation language, booleans are a
subtype of integers: \texttt{True == 1} evaluates true, the two hash
identically, and a dictionary keyed by \texttt{1} silently swallows a later
entry keyed by \texttt{True}. A finite transition table keyed by untagged
atoms could therefore conflate a numeric state $1$ with a logical state
\emph{true} --- merging two rows of the table, changing its size, and
corrupting its certificate --- without any error being raised. The tagged
payload assigns the two atoms different type tags, hence different
canonical texts, hence different digests; the kernel-facing
\texttt{TaggedFiniteValue} token exists precisely so that hash-based
containers ``whose native equality conflates types'' (the implementation's
own words) can key on the tag pair instead of the raw atom.

\emph{The string-coercion pun.} Configuration formats, spreadsheets, and
lenient parsers routinely deliver \texttt{"true"} where a boolean was
meant, or \texttt{"1"} where an integer was meant. A permissive layer that
coerced these would make a board's semantics depend on which serialization
path a parameter happened to travel. In the candidate, a string parameter
is a string parameter: the tagged identity of \texttt{"true"} shares
nothing with the tagged identity of \emph{true}, canonical validation
rejects a payload whose tag and value disagree (an \texttt{int} tag whose
value string is not the exact canonical decimal form; a \texttt{bool} tag
whose value is not a JSON boolean), and the round-trip check of
Section~\ref{ch04:sec:traceid} ensures the distinction survives
serialization.

\emph{The evidence pun.} The deepest reason is evidential. Identities are
hashes of canonical bytes; if two semantically distinct values could share
canonical bytes, the identity system would licence false equalities in the
permanent record. Tagging makes the value space injective into the payload
space by construction: distinct (tag, value) pairs yield distinct
canonical texts. Where floating-point values must enter a canonical
payload (board metadata admits them), they are serialized by exact hexadecimal
representation with explicit \texttt{NaN} and infinity markers --- again an
injective encoding, chosen so that ``same digest'' can never paper over
``different value.''

\section{The single immutable decimal context}\label{ch04:sec:decimal}

All decimal arithmetic in the candidate runs under one versioned context
profile, \texttt{cas-pcba-v8.decimal-context.v1}: precision $60$;
rounding half-even; exponent range $E_{\min} = -999{,}999$ to
$E_{\max} = 999{,}999$; clamping off; capitals on; and traps armed for
division by zero, invalid operation, and overflow.

\begin{ruleenv}[Per-operation context creation]\label{ch04:rule:context}
The context is created \emph{from primitive constants} for each public
operation, used for exactly that operation's dynamic extent, and discarded.
No caller-supplied or ambient (thread-local) decimal context is ever
consulted; a caller running under precision $3$, or under a context with
traps disarmed, obtains bit-identical results to a caller running under
defaults. Public arithmetic entry points are wrapped so that the fixed
context is installed before the operation begins and uninstalled after it
completes.
\end{ruleenv}

The rule has an evidence corollary. The complete context payload is hashed
once, at import, into the constant \texttt{DECIMAL\_CONTEXT\_SHA256}, and
that digest enters: every decimal value identity in a trace; every
mixed-simplex identity payload; every projector identity payload; every
component identity payload; every board identity payload; and the engine
identity itself. A future reader of any artifact can therefore verify not
only \emph{what} was computed but \emph{under which arithmetic} --- and any
change to the arithmetic profile, however small, changes every downstream
digest, making silent arithmetic drift structurally impossible.

Three trapped conditions replace three silent failure modes: division by
zero raises instead of yielding an infinity; invalid operations raise
instead of yielding \texttt{NaN}; overflow raises instead of saturating.
Since the type layer separately excludes non-finite decimals from every
record field (weights, uncertainties, bounds, gains must be finite), the
combination gives a closed world: no \texttt{NaN} or infinity can be
constructed by candidate arithmetic, and none can be smuggled in through a
field.

Alongside context-governed arithmetic, the candidate provides
\emph{context-independent} exact summation: finite decimals are summed by
aligning exponents and adding integer coefficients, with no rounding at any
precision. This is the operation behind the exact sum-to-one constraint of
Definition~\ref{ch04:def:simplex} and the exact residual computation of
Section~\ref{ch04:sec:normalization}: a constraint checked by exact
summation is an identity of integers, not a tolerance judgment.

\begin{remark}[Relation to IEEE 754]\label{ch04:rem:ieee}
IEEE~754 \cite{ieee754} is cited here as the external reference point for
explicitly controlled computer arithmetic: reproducible results require a
fully specified format, rounding rule, and exception discipline, and the
candidate adopts exactly that posture. But the candidate's decimal profile
is its \emph{own} software contract, not an instance of an IEEE
interchange format: precision $60$ is not an IEEE decimal format
parameter, the exponent range is the candidate's, and conformance claims
to the standard are neither made nor needed. The standard supplies the
engineering lesson (specify the arithmetic, or lose determinism); the
profile supplies the specification; the digest supplies the proof of which
specification was in force. Where binary64 floating point \emph{is} used
--- solely inside the disclosed Fisher--Rao pseudometric of
Chapter~\ref{ch06:chapter} --- that use is named, profiled, and bounded,
and its limits are part of the disclosure rather than an ambient
assumption.
\end{remark}

\section{Canonical identity is injective}\label{ch04:sec:injective}

The tagged-value discipline of Section~\ref{ch04:sec:tagged} and the
canonical serializations of Section~\ref{ch04:sec:traceid} exist to
support one property, which is worth stating and proving at the level of
the declared encodings.

\begin{proposition}[Injectivity of the finite-atom encoding]
\label{ch04:prop:injective}
Let $V = \mathbb{B} \sqcup \mathbb{Z} \sqcup \Sigma^{*}$ be the disjoint
union of booleans, integers, and finite strings, and let
$e : V \to \Sigma^{*}$ be the canonical tagged encoding: $e(v)$ is the
canonical JSON text of the payload carrying the profile constant, the type
tag of the summand $v$ inhabits, and the canonical value representation
(JSON boolean; exact decimal integer string; the string itself). Then $e$
is injective: $e(v) = e(w)$ implies $v = w$, \emph{including} across
summands.
\end{proposition}

\begin{proof}
Suppose $e(v) = e(w)$. Canonical JSON serialization with sorted keys and
fixed separators is injective on payload objects, so the payloads of $v$
and $w$ are equal as objects; in particular their type tags agree, so $v$
and $w$ inhabit the same summand. Within a summand, the value
representation is injective: JSON booleans distinguish the two booleans;
the exact decimal string of an integer determines the integer (the
canonical form forbids leading zeros and non-canonical signs, and the
validator rejects any value string that does not round-trip through
integer parsing to itself); a string represents itself. Hence $v = w$.
The cross-summand cases that defeat untagged encodings --- integer $1$
versus Boolean \emph{true} versus string \texttt{"true"}, or integer $1$
versus string \texttt{"1"} --- are separated at the first step, by the
type tag.
\end{proof}

\begin{corollary}[Digest separation for finite atoms]
\label{ch04:cor:digest}
Distinct finite atoms have distinct canonical texts, hence distinct
SHA-256 identities up to hash collision. The candidate treats SHA-256
collision resistance as an external assumption, disclosed as such: no
claim in this monograph depends on it beyond the standard evidentiary use
of content digests.
\end{corollary}

The corollary marks a contrast that Chapter~\ref{ch06:chapter} develops at
length: the \emph{identity} layer of the candidate separates all distinct
values by construction (modulo the disclosed hash assumption), while the
\emph{numerical geometry} layer explicitly does not promise separation at
binary64 resolution. Claims of sameness belong to the identity layer;
claims of closeness belong to the geometry layer; the type system keeps
the two vocabularies apart so that neither can borrow the other's
strength.

\section{Measurement kinds as a claim taxonomy}\label{ch04:sec:measkinds}

The six-element \texttt{MeasurementKind} enumeration of
Definition~\ref{ch04:def:calref} is small, but it is the type-level
skeleton of the candidate's entire evidence discipline, so its semantics
deserve explicit statement.

\begin{description}
\item[\textsc{design\_value}.] A quantity asserted by the board author as
part of the design: a parameter, a fixture input, a declared constant. It
carries no measurement claim at all; it is structural metadata. The
default measurement kind of every input sample is \textsc{design\_value},
so a fixture that says nothing claims nothing.
\item[\textsc{behavioral\_proxy}.] A quantity derived from observed
behavior of some system under a declared protocol, offered as a proxy for
a modeled quantity. The kind records that the link from observation to
quantity is itself a modeling assumption.
\item[\textsc{instrumented}.] A quantity attributed to an instrument
chain with declared identity, protocol, and uncertainty. In the v0.1
candidate this kind can be \emph{declared} but never \emph{authenticated}:
no instrument chain exists in the admitted evidence, and locally declared
instrumented records remain structural/test-only.
\item[\textsc{human\_score}, \textsc{llm\_score}.] Quantities produced by
a human rater or a language-model rater under a declared protocol. The
kinds exist so that rater-derived numbers can never be conflated with
instrumented ones by type.
\item[\textsc{model\_derived}.] A quantity computed by the candidate
itself from other quantities. Board \emph{inputs} may not carry this kind
(Section~\ref{ch04:sec:tpcal}); it exists for derived, child measurements
inside a trace, where the parent linkage documents the derivation.
\end{description}

The enumeration is thus a closed taxonomy of \emph{how a number came to
be}, attached at the type level to every number that enters the evidence.
The DRC and simulator consume it (calibration compatibility, input
admissibility); the Evidence-Packet discipline consumes it (eligibility by
kind); and the audit consumes it (a report can be filtered, by type, to
exactly the claims of a given evidentiary strength).

\section{Worked type-level fragments}\label{ch04:sec:worked}

The fragments below are \emph{illustrative only}: they are not normative
source code, they elide imports and error text, and they exist to make the
failure modes of this chapter concrete. Chapter~\ref{ch05:chapter} gives
the corresponding board-level constructions.

\begin{example}[Legal state literals]\label{ch04:ex:legal}
The partial literal $({\wedge}3, P, \bot, \bot)$ --- constructed as
\texttt{CasState(noetic=Noetic.NEGATIVE, domain=Domain.PRINCIPLED)} ---
is legal, renders as \texttt{\^{}3\^{}P}, and projects to the cell
\texttt{P3} (index $4$ in Table~\ref{ch04:tab:cells}). The full literal
\texttt{CasState(Noetic.NEGATIVE, Domain.PRINCIPLED, Drive.F2,
Register.ANALYTICAL)} is equally legal, renders as
\texttt{\^{}3\^{}P\_\{F2\}|A}, and projects to the \emph{same} cell ---
the worked instance of Proposition~\ref{ch04:prop:manytoone}.
\end{example}

\begin{example}[Illegal state literals]\label{ch04:ex:illegal}
Each of the following is rejected at construction with a typed error, not
repaired: \texttt{CasState(3, Domain.PRINCIPLED)} (bare integer where a
\texttt{Noetic} is required); \texttt{CasState(Noetic.NEGATIVE, "P")}
(bare string where a \texttt{Domain} is required);
\texttt{CasState(Noetic.NEGATIVE, Register.EMBODIED)} (a Register in the
Domain position --- rejected even though both carry the string \texttt{E}
family, because the enumerations are disjoint types); and any instance of
a subclass of \texttt{CasState}, however faithful its fields
(Rule~\ref{ch04:rule:exact}).
\end{example}

\begin{example}[Simplex boundary rejection]\label{ch04:ex:simplexreject}
The uniform profile with all forty weights \texttt{Decimal("0.025")} is a
legal address: weights are strictly positive and their exact sum is
\texttt{Decimal("1")}. The ``one-hot'' tuple with weight \texttt{1} on
\texttt{P0} and \texttt{0} elsewhere is rejected (zero coordinates); the
tuple of forty weights \texttt{0.02} and \texttt{0.03} alternating is
rejected if its exact sum differs from one; a tuple of thirty-nine
weights is rejected on arity; and a tuple containing
\texttt{float 0.025} rather than \texttt{Decimal("0.025")} is rejected on
value type, because binary64 atoms have no place in an exact address.
\end{example}

\begin{example}[Type-pun rejection]\label{ch04:ex:pun}
A finite transition table over carrier $\{0, 1\}$ and a table over carrier
$\{\mathit{false}, \mathit{true}\}$ have disjoint tagged carriers and
distinct certificates, even though the underlying language would happily
evaluate \texttt{True == 1} as true. A fixture that attempts to feed the
string \texttt{"true"} to a decision port fails port validation: the
serialized identity has tag \texttt{string}, and a decision port requires
tag \texttt{bool} (Section~\ref{ch04:sec:traceid}).
\end{example}

\begin{example}[Ambient-context irrelevance]\label{ch04:ex:ambient}
A caller that sets its thread-local decimal precision to $3$ and then
normalizes the weight tuple $(1, 1, 1, \ldots, 1)$ (forty ones) obtains
the same address, with the same forty weights of
\texttt{Decimal("0.025")} and the same profile digest, as a caller running
under default precision: the normalization runs under the fixed
sixty-digit context of Rule~\ref{ch04:rule:context} regardless of ambient
state, and the exact-sum check is context-free integer arithmetic. The
ambient precision is not consulted, not restored incorrectly, and not
recorded --- because it never mattered.
\end{example}

\section{Invariant summary}\label{ch04:sec:invariants}

\begin{invariant}[Type-layer invariants]\label{ch04:inv:summary}
For every value admitted by the type system:
(i) its runtime type is exact --- no subclass, no coercion;
(ii) every factor field is an exact enumeration member, with absence
explicit where permitted and nowhere defaulted;
(iii) every reference string satisfies the safe-reference profile (with
the single audited \texttt{\$board} exception);
(iv) every decimal field is finite, and all decimal arithmetic occurred
under the fixed versioned context, whose digest accompanies the value;
(v) every finite atom carries its type tag, and no two distinct tagged
atoms share a canonical identity;
(vi) every simplex address has exactly forty strictly positive weights in
canonical cell order with exact sum one;
(vii) every evidence-bearing record binds the fixed research boundary and
is identified by the SHA-256 of its complete canonical payload; and
(viii) revalidation of any record at any later time either reconfirms all
of the above or raises --- there is no degraded acceptance path.
\end{invariant}

\begin{threat}[What the type system does not do]\label{ch04:threat:limits}
The type system validates form, not truth. A well-typed calibration
reference does not evidence calibration; a well-typed state binding does
not evidence that any system occupies that state; a well-typed simplex
address does not evidence that forty numbers describe anything outside the
software. Type-level guarantees are also bounded by the implementation's
own correctness: the invariants above are enforced by tested code, not
proved code, and the bounded-evidence posture of
Chapter~12 applies to them in full.
\end{threat}

\section{Claim boundary}\label{ch04:sec:boundary}

\begin{boundary}\label{ch04:bnd:main}
This chapter specifies software contracts of the CAS PCBA~8 v0.1 research
candidate, at Ring~4, \textsc{not licensed}. The types named after CAS
coordinates represent those coordinates inside the candidate only; no
claim is made that any human, biological, or artificial system realizes
them. Strict positivity, exact summation, tagged identity, and the fixed
decimal context are properties of the candidate's arithmetic and records,
not of any external quantity. No physical, causal, safety, clinical,
regulatory, or product claim is made or licensed by anything in this
chapter, and the projection, positivity, and identity results stated here
(Propositions and Rules \ref{ch04:prop:manytoone},
\ref{ch04:rule:notonehot}, \ref{ch04:rule:nottks},
\ref{ch04:rule:typepun}, \ref{ch04:rule:context}) are statements about
declared types and tested code paths, offered for independent review.
\end{boundary}
